\scanpage{99}
(sur une extension $k'$ de $k$ contenant $\unsure{\zeta|\rho},\zeta'$) sont les droites homogènes
\[
\begin{aligned}
&\deleted{\unsure{Y/X|X}=\zeta\ \text{et}}\\
&Y-\zeta X=0\qquad Y-\zeta'X=0
\end{aligned}
\]
qui sont aussi les droites propres sous $u_0$
\[
(\qquad u_0(X,Y)=(-Y+\alpha X,X)
\]
[ce qui correspond au fait que $u_0$ est une \deleted{\unsure{(pour|pour} \illegible{1}} \uline{\unsure{rotation|rotatation}} pour $f_0$, si $f_0$ est non dég. i.e. $\alpha\ne\pm2$]

Dans le cas $\alpha=\pm2$, ces deux directions sont confondues ; \deleted{\unsure{si}} Pour $\alpha=-2$, $\zeta=\zeta'=-1$, on trouve la direction
\[
X\unsure{-|+}Y=0
\]
qui apparaissait déjà comme privilégiée dans l'étude du ``cas II'' \deleted{($\alpha=-2$)}. Pour $\alpha=2$, $\zeta\unsure{=}\zeta'=1$, c'est la direction $X=Y$, \deleted{\unsure{qu}} on va l'examiner plus en détail plus bas.

8) Discriminant \inserted{divisé} de $f$. La matrice de $f$ est
\[
\begin{pmatrix}2&-\alpha&-1\\-\alpha&2&-1\\-1&-1&\deleted{\unsure{2}}0\end{pmatrix}.
\]
son discriminant est égal à
\[
\delta=2\delta'\qquad\delta'=-(\alpha+2)
\]
% archive sheet number: 98
